\documentclass[11pt]{article}
\usepackage[margin=1in]{geometry}
\title{Paper Moons of Caldera}
\author{Fictional Movie Catalog}
\date{2026}

\begin{document}
\maketitle

\section*{Catalog Metadata}
\textbf{Release date:} November 6, 2026\\
\textbf{Runtime:} 102 minutes\\
\textbf{Genres:} Animation, Family, Science Fiction\\
\textbf{Director:} Lumi Aster\\
\textbf{Writers:} Priya Noll and Evan Moss\\
\textbf{Top cast:} Aya Lin as Niko; Marco Day as Dr. Sol; June Kestrel as Pip;
Rafi Bloom as Commander Vale.\\
\textbf{Mock IMDb rating:} 8.0 from 12,780 votes

\section*{Synopsis}
Twelve-year-old Niko lives in Caldera Station, a research colony suspended
above a volcanic planet. She folds paper models of the planet's tiny moons and
maps their paths across her bedroom ceiling. When her mother's survey ship
loses navigation during an ash eruption, official instruments cannot see
through the charged clouds.

Niko notices that her paper moons drift whenever magnetic waves pass through
the station. With help from the maintenance robot Pip, she builds a fleet of
lightweight signal gliders based on her models. The gliders form a moving
constellation that the lost crew can follow home. Commander Vale initially
dismisses the plan, but Niko's careful observations reveal a safe route no
computer model has considered.

\section*{Themes and Style}
The film celebrates curiosity, craft, and young people whose knowledge is
overlooked because it arrives in an unexpected form. Hand-painted paper
textures appear throughout the computer animation. The volcanic clouds shift
from threatening charcoal shapes to a luminous night sky as Niko learns to
trust her own method.

\section*{Production Notes}
Artists built physical origami models for every spacecraft before animation
began. The score blends a children's choir with sounds recorded from folding
paper. The fictional production is rated PG for mild peril and emotional
themes.

\end{document}
